
%\subsection*{Proof of Theorem~\ref{thm-generalization-bound}}
\subsection*{A4. Proof of Theorem~\ref{thm-generalization-bound}}
We establish the Rademacher complexity  of \HFrame
by extending the proof of 
the Rademacher complexity for \GNNs~\cite{Gen-bound}.
Since \dual does not have any parameters, 
and works for any inputs, we only 
bound the Rademacher complexity of \HGIN.
We prove the bound in multiple steps:
(1) represent the computation of \HGIN
by a tree;
(2) quantify the impact of change 
of parameters on the embeddings;
(3) bound the changes in prediction probability;
(4) bound the covering numbers for \HGIN;
and (5) finally bound the Rademacher complexity
using the covering numbers.

In the following, we
show these steps one by one. 
In the proof of Theorem~\ref{thm-upper-bound}, 
the computation of \HGIN is represented by a 
computation tree.
It remains to prove steps (2)-(5).	
	
\stitle{$\bullet$ Step 2: Bound the impact of parameters}.
We first construct a formula 
to represent changes of embeddings, and then bound the changes.


Based on Equations (1)-(4), 
we can deduce the following upper bound
for the changes of embeddings:
\begin{small}
\begin{align*}
\Delta_{L}  \triangleq & \quad\|T_{L}(W_{1},\mathds{W})-T_{L}(W_{1}^{\prime},\mathds{W}^{\prime})\|_{2} \\
  =&\quad\bigg\|\kw{COM}(W_{1}x_{L} + \sum_{r\in R} W_r\underbrace{\rho(\sum_{j\in  N^+_r(x_{L})}\kw{MSG}_{{\bf 1}[j=v]}(T_{L-1,j}(W_{1},\mathds{W}))}_{\overset{\Delta}
 {\operatorname*{\operatorname*{=}}}R^+(W_{1},\mathds{W},x_{L},r)}),\\
 &\quad  W_{1}x_{L}+\sum_{r\in R} W_r\underbrace{\rho(\sum_{j\in 
 		N^-_r(x_{L})}\kw{MSG}_{{\bf 1}[j=v]}(T_{L-1,j}(W_{1},\mathds{W}))}_{\overset{\Delta}
 	{\operatorname*{\operatorname*{=}}}R^-(W_{1},\mathds{W},x_{L},r)})\\
 & -\quad\kw{COM}(W'_{1}x_{L}+\sum_{r\in R} W'_r\rho(\sum_{j\in 
 		N^+_r(x_{L})}\kw{MSG}_{{\bf 1}[j=v]}(T_{L-1,j}(W'_{1},\mathds{W}'))), \\
 		&\quad  W'_{1}x_{L}+\sum_{r\in R} W'_r\rho(\sum_{j\in 
 		N^-_r(x_{L})}\kw{MSG}_{{\bf 1}[j=v]}(T_{L-1,j}(W'_{1},\mathds{W}')))\bigg\|_{2} \\
  \leq& \quad 2C_{\kw{COM}}\left\|(W_{1}-W_{1}^{\prime})x_{L}\right\|_{2} \\
 & + C_{\kw{COM}}\sum_{r\in R}\|(W_rR^+(W_{1},\mathds{W},x_{L})-W_r^{\prime}R^+(W_{1}^{\prime},\mathds{W}^{\prime},x_{L}))\|_{2}\\
 & +\quad C_{\kw{COM}}\sum_{r\in R}\|(W_rR^-(W_{1},\mathds{W},x_{L})-W_r^{\prime}R^-(W_{1}^{\prime},\mathds{W}^{\prime},x_{L}))\|_{2}  \tag{10}\\
  \leq& \quad 2C_{\kw{COM}}\left\|(W_{1}-W_{1}^{\prime})x_{L}\right\|_{2} \\
& +\quad C_{\kw{COM}}\sum_{r\in R}\|(W_rR^+(W_{1},\mathds{W},x_{L})-W_r^{\prime}R^+(W_{1},\mathds{W},x_{L}))\|_{2}\\
& + C_{\kw{COM}}\sum_{r\in R}\|(W^{\prime}_rR^+(W_{1},\mathds{W},x_{L})-W_r^{\prime}R^+(W_{1}^{\prime},\mathds{W}^{\prime},x_{L}))\|_{2}\\
& +\quad C_{\kw{COM}}\sum_{r\in R}\|(W_rR^-(W_{1},\mathds{W},x_{L})-W_r^{\prime}R^-(W_{1},\mathds{W},x_{L}))\|_{2}\\
& + C_{\kw{COM}}\sum_{r\in R}\|(W^{\prime}_rR^-(W_{1},\mathds{W},x_{L})-W_r^{\prime}R^-(W_{1}^{\prime},\mathds{W}^{\prime},x_{L}))\|_{2} \tag{11}
\end{align*}
\end{small}

Equation~(10) holds, since we assume that 
 function \kw{COM} is $C_{\kw{COM}}$-Lipschitz.
 
 


To bound $\Delta_{L}$, 
we bound the following:
\begin{itemize}
\item[(a)] $\|(W_{1}{-}W_{1}^{\prime})x_{L}\|_{2}$;
\item[(b)]  $\|W_rR^+(W_{1},\mathds{W},x_{L}){-}W_r^{\prime}R^+(W_{1},\mathds{W},x_{L})\|_{2}$;
\item[(c)]  $\|W^{\prime}_rR^+(W_{1},\mathds{W},x_{L}){-}W_r^{\prime}R^+(W_{1}^{\prime},\mathds{W}^{\prime},x_{L})\|_{2}$;
\item[(d)]  $\|W_rR^-(W_{1},\mathds{W},x_{L})-W_r^{\prime}R^-(W_{1},\mathds{W},x_{L})\|_{2}$;
and 
\item[(e)]  $\|(W^{\prime}_rR^-(W_{1},\mathds{W},x_{L}){-}W_r^{\prime}R^-(W_{1}^{\prime},\mathds{W}^{\prime},x_{L})\|_{2}$.
\end{itemize}
%\looseness=-1

For term (a), we can assume an upper bound
for the norm of $W_1$ and $x_L$. 
For (b), 
we have 
$\|W_rR^+(W_{1},\mathds{W},x_{L})-W_r^{\prime}R^+(W_{1},\mathds{W},x_{L})\|_{2}
\leq \|(W_r-W_r^{\prime})R^+(W_{1},\mathds{W},x_{L})\|_{2}$.
Then we only need to bound $\|R^+(W_{1}^{\prime},\mathds{W}^{\prime},x_{L})\|_{2}$.
For (c), 
we have that 
$\|W^{\prime}_rR^+(W_{1},\mathds{W},x_{L})-W_r^{\prime}R^+(W_{1}^{\prime},\mathds{W}^{\prime},x_{L})\|_{2}
\leq \|W^{\prime}_r(R^+(W_{1},\mathds{W},x_{L})-R^+(W_{1}^{\prime},\mathds{W}^{\prime},x_{L}))\|_{2}$.
Therefore, we bound 
$\|R^+(W_{1},\mathds{W},x_{L})-R^+(W_{1}^{\prime},\mathds{W}^{\prime},x_{L})\|_{2}$.
The cases (d) and (e) can be analyzed similarly.
In the following, we only provide
upper bounds for $\|R^+(W_{1}^{\prime},\mathds{W}^{\prime},x_{L})\|_{2}$
and $\|R^+(W_{1},\mathds{W},x_{L})-R^+(W_{1}^{\prime},\mathds{W}^{\prime},x_{L})\|_{2}$.

\vspace{2ex}

\eat{
analyze the terms containing $R^+(W_{1},\mathds{W},x_{L})$.
The terms containing $R^+(W_{1},\mathds{W},x_{L})$ 
can be analyzed similarly.
}


\etitle{Upper bound for $\|R^+(W_{1},\mathds{W},x_{L})-R^+(W_{1}^{\prime},\mathds{W}^{\prime},x_{L}))\|_{2}$}.
We start with the upper bounds for term $\|R^+(W_{1},\mathds{W},x_{L})-R^+(W_{1}^{\prime},\mathds{W}^{\prime},x_{L}))\|_{2}$.
Based on the definition of 
$R^+(W_{1},\mathds{W},x_{L})$, we have the following.

\eat{
\begin{equation*}\begin{aligned}
&\|W_r^{\prime}R^+(W_1,\mathds{W},x_L)-W_r^{\prime}R^+(W_1^{\prime},\mathds{W}^{\prime},x_L)\|_2 
\leq&\|W_r^{\prime}\|_2\|R^+(W_1,\mathds{W},x_L)-R^+(W_1^{\prime},\mathds{W}^{\prime},x_L)\|_2
\end{aligned}\end{equation*}
}


\begin{scriptsize}	
%\vspace{-7ex}
\begin{align*}
&\|R^+(W_1,\mathds{W},x_L,r)-R^+(W_1',\mathds{W}',x_L,r)\|_2\\ % \tag{eq-1} 
  &\leq C_\rho\bigg\|\sum_{j\in N_r^+(x_L)}\bigg(\kw{MSG}_{{\bf 1}[j=v]}(T_{L-1,j}(W_1,
  \mathds{W})) -\kw{MSG}_{{\bf 1}[j=v]}(T_{L-1,j}(W_1',\mathds{W}'))\bigg)\bigg\|_2  \tag{12}
  \\
 & \leq C_\rho\sum_{j\in N_r^+(x_L)}\bigg\|\bigg(\kw{MSG}_{{\bf 1}[j=v]}(T_{L-1,j}(W_1,
  \mathds{W})) -\kw{MSG}_{{\bf 1}[j=v]}(T_{L-1,j}(W_1',\mathds{W}'))\bigg)\bigg\|_2 \nonumber
  \\
  &\leq C_\rho C_{\kw{MSG}_0}C_{\kw{MSG}_1}\sum_{j\in N_r^+(x_L)}\bigg\|T_{L-1,j}(W_1,\mathds{W})-
  T_{L-1,j}(W_1',\mathds{W}')\bigg\|_2 \tag{13}\\
  &= C_\rho C_{\kw{MSG}_0}C_{\kw{MSG}_1}\sum_{j\in N_r^+(x_L)}\Delta_{L-1,j}\\
  &\leq C_\rho C_{\kw{MSG}_0}C_{\kw{MSG}_1}d\max_{j\in N_r^+(x_L)}\Delta_{L-1,j}\tag{14}
\end{align*}
\end{scriptsize}


Equations (12) and (13) follow from the assumption 
that Functions $\rho$, $\kw{MSG}^{k,r}_0$ and 
$\kw{MSG}^{k,r}_0$ are $C_\rho$-Lipschitz, 
$C_{\kw{MSG}_0}$-Lipschitz  and $C_{\kw{MSG}_1}$-Lipschitz,
respectively.
Equation (14) holds, since we assume that 
the degrees of vertices in $G$ are bounded by 
a constant $d$.
\looseness = -1

\eat{
Since we assume that $\|W'_r\|\leq B_2$,
we can deduce that the bound as follows.

\begin{equation*}\begin{aligned}
||W_r^{\prime}R(W_1,\mathds{W},x_L,r)  - W_r^{\prime}R(W_1^{\prime},\mathds{W}^{\prime},x_L,r)\|_2
 &\leq B_2C_\rho C_{\kw{MSG}_0}C_{\kw{MSG}_1}\sum_{j\in N_r^+(x_L)}\Delta_{L-1,j} \\
 &\leq B_2C_\rho C_{\kw{MSG}_0}C_{\kw{MSG}_1}d\max_{j\in N_r^+(x_L)}\Delta_{L-1,j}
\end{aligned}\end{equation*}


\begin{equation*}\begin{aligned}
\Delta_{L} & \leq\quad 2C_{\kw{COM}} B_x\left\|(W_1-W_1^{\prime})\right\|_2 
  +\quad C_{\kw{COM}} |R| B_2C_\rho C_{\kw{MSG}_0}C_{\kw{MSG}_1}d\max_{j\in N_r^+(x_L)}\Delta_{L-1,j} \\
 &\quad +\kw{COM}|R|\max_{r\in R}\|(W_r-W_r^{\prime})R^+(W_1,\mathds{W},x_L,r)||_2
  + \kw{COM}|R|\max_{r\in R}\|(W_r-W_r^{\prime})R^-(W_1,\mathds{W},x_L,r)||_2
\end{aligned}\end{equation*}
}

Observe that the upper bound for
$\Delta_{L}$ depends on that of $\|R^+(W_1,\mathds{W},x_L,r)-R^+(W_1',\mathds{W}',x_L,r)\|_2$,
which in turn dependents on $\Delta_{L-1,j}$.
Then we deduce a bound for $\|R^+(W_1,\mathds{W},x_L,r)-R^+(W_1',\mathds{W}',x_L,r)\|_2$
by expanding the bounds (see below). 


% \begin{equation}
% \begin{aligned}
% \|R(W_1,\mathds{W},x_L)\|_2 \\ 
% \leq\quad C_\rho C_gd\min\left\{b\sqrt{r},C_\phi B_1B_x\frac{(\mathcal{C}d)^L-1}{\mathcal{C}d-1}\right\}
% \end{aligned}
% \end{equation}



\etitle{Upper bounds for $\|R^+(W_{1}^{\prime},\mathds{W}^{\prime},x_{L}))\|_{2}$}.
Based on the definition of $\|R^+(W_{1},\mathds{W},x_{L},r)\|_{2}$, 
we can deduce the following.

\begin{align*}
  \|R^+(W_{1},\mathds{W},x_{L},r)\|_{2} 
 & =\bigg\|\rho\big(\sum_{j\in N_r^+(x_{L})}\kw{MSG}_{{\bf 1}[j=v]}(T_{L-1,j}(W_{1},\mathds{W})
 \big)\big)\bigg\|_{2} \\
 & \leq C_\rho\bigg\|\sum_{j\in N_r^+(x_L)}\kw{MSG}_{{\bf 1}[j=v]}(T_{L-1,j}(W_1,\mathds{W}))\bigg\|
 _2 \\
 & \leq C_\rho\sum_{j\in N_r^+(x_L)}\bigg\|\kw{MSG}_{{\bf 1}[j=v]}(T_{L-1,j}(W_1,
 \mathds{W}))\\
 &\quad\quad -\kw{MSG}_{{\bf 1}[j=v]}(0)\bigg\|_2 \\
 & \leq C_{\rho}C_{\kw{MSG}_0}C_{\kw{MSG}_1}\sum_{j\in N_r^+(x_{L})}\bigg\|T_{L-1,j}(W_{1},\mathds{W})
 \bigg\|_{2} \\
 & \leq C_\rho C_{\kw{MSG}_0}C_{\kw{MSG}_1}d\max_{j\in N_r^+(x_L)}\bigg\|T_{L-1,j}(W_1,\mathds{W})\bigg\|_2
\end{align*}


The equations hold due to the reasons given above.
Observe that $\|R^+(W_{1},\mathds{W},x_{L},r)\|_{2}$ 
is bounded using $\bigg\|T_{L-1,j}(W_1,\mathds{W})\bigg\|_2$,
\ie the changes of the predictions of neighbors.
We continue to expand $\bigg\|T_{L-1,j}(W_1,\mathds{W})\bigg\|_2$
as follows.


\begin{align*}
		&\quad\quad\bigg\|T_{L-1,j}(W_1,\mathds{W})\bigg\|_2\\
		 &=\quad\bigg\|\kw{COM}(W_{1}x_{L}+\sum_{r\in R} W_r R^+(W_1,W_2,x_{L-1,j},r),
		\\
		&\quad\quad\quad W_{1}x_{L}+\sum_{r\in R} W_r R^-(W_1,W_2,x_{L-1,j},r)) \bigg\|
_2 \\
		& =\quad\bigg\|\kw{COM}(W_{1}x_{L}+\sum_{r\in R} W_r R^+(W_1,W_2,x_{L-1,j},r),
W_{1}x_{L}\\
&\quad\quad\quad+\sum_{r\in R} W_r R^-(W_1,W_2,x_{L-1,j},r)) -
\kw{COM}(0,0) \bigg\|
_2 \\
		& \leq\quad C_{\kw{COM}}\bigg\|W_{1}x_{L-1,j}+\sum_{r\in R} W_r 
		R^+(W_{1},W_{2},x_{L-1,j},r)\bigg\|_{2} \\
		&\quad\quad+  C_{\kw{COM}}\bigg\|W_{1}x_{L-1,j}+\sum_{r\in R} W_r 
R^-(W_{1},W_{2},x_{L-1,j},r)\bigg\|_{2} \tag{15} \\
		& \leq\quad 2C_{\kw{COM}}\bigg\|W_{1}x_{L-1,j}\bigg\|_{2}+C_{\kw{COM}}\sum_{r\in R} \bigg\|
		W_rR^+(W_{1},W_{2},x_{L-1,j},r)\bigg\|_{2} \\
		&\quad\quad+ C_{\kw{COM}}\sum_{r\in R} \bigg\|
		W_rR^-(W_{1},W_{2},x_{L-1,j},r)\bigg\|_{2}\\
		& \leq\quad 2C_{\kw{COM}}B_{1}B_{x}+C_{\kw{COM}}|R|B_{2}\bigg\|
		R^+(W_{1},W_{2},x_{L-1,j},r)\bigg\|_{2}\\
		&\quad\quad+C_{\kw{COM}}|R|B_{2}\bigg\|
		R^-(W_{1},W_{2},x_{L-1,j},r)\bigg\|_{2}.  \tag{16}
\end{align*}


Equation (15) holds, since $\kw{COM}$
is $C_{\kw{COM}}$-Lipschitz.
And Equation (16) holds, since 
we assume that 
the norms of $W_1$ and $X_i$
is bounded by $B_1$ and $B_x$, respectively.
%
Using the above bound, we bound $\|R^+(W_1,\mathds{W},x_L,r)\|_2$ %can be bounded
as follows.

\begin{align*}
&  \|R^+(W_1,\mathds{W},x_L,r)\|_2 \\
 & \leq 2C_\rho C_{\kw{MSG}_0}C_{\kw{MSG}_1}C_{\kw{COM}} B_1B_xd  \\
 & +  C_\rho C_{\kw{MSG}_0}C_{\kw{MSG}_1}C_{\kw{COM}} |R| 
 B_2d\max_{j\in C(x_L)}\bigg(\bigg\|R^+(W_1,
 \mathds{W},x_{L-1,j})\bigg\|_2\\
 &\quad\quad+\bigg\|R^-(W_1,
 \mathds{W},x_{L-1,j})\bigg\|_2\bigg) \\
 & \leq 2C_\rho C_{\kw{MSG}_0}C_{\kw{MSG}_1}C_\phi B_1B_xd\sum_{\ell=0}^{L-1}(2C_\rho C_{\kw{MSG}_0}C_{\kw{MSG}_1}
 C_{\kw{COM}} B_2 |R| d)^\ell \\
 & =\quad 2C_\rho C_{\kw{MSG}_0}C_{\kw{MSG}_1}C_{\kw{COM}} B_1B_xd\frac{(\mathcal{C}d)^L-1}{\mathcal{C}d-1}. \tag{17}
\end{align*}

Let ${\mathcal C} = 2C_\rho C_{\kw{MSG}_0}C_{\kw{MSG}_1}C_{\kw{COM}} B_2 |R| d$.
Denote $M$ as $\frac{\left(\mathcal{C}d\right)^{L}-1}{\mathcal{C}d-1}$.
That is, $\|R^+(W_1,\mathds{W},x_L,r)\|_2  \leq {\mathcal C} B_1  B_xM$.

\eat{
Moreover, since we assume that $\|x\|_\infty=|x_1|+\ldots+|x_r|\leq 
b$ for each vector $x=(x_1, \ldots, x_r)\in R^r$, 
$\|x\|_2\leq \sqrt{x^2_1+\ldots + x^2_r} \leq \sqrt{r}\|x\|_\infty$
by the mean value inequality.
We have that 
$$\left\|T_{L-1,j}(W_1,\mathds{W})\right\|_2\leq b\sqrt{r}$$

Moreover, we have that 

\begin{align*}\|R^+(W_1,\mathds{W},x_L,r)\|_2\leq bdC_\rho C_{\kw{MSG}_0}C_{\kw{MSG}_1}\sqrt{r}
\end{align*}


Combining \tbf and \tbf, we have that
\begin{align*}\|R^+(W_1,\mathds{W},x_L,r)\|_2\leq C_\rho C_{\kw{MSG}_0}C_{\kw{MSG}_1}d\min\left\{b\sqrt{r},C_\kw{COM} 
B_1B_x\frac{(\mathcal{C}d)^L-1}{\mathcal{C}d-1}\right\}\end{align*}
}


\etitle{Upper bounds for $\Delta_{L}$}.
Combining Equations (11), (14) and (17), 
we get

\eat{
\begin{gather*}
\Delta_{L}\quad\leq\quad MB_{x}\left\|W_{1}-W_{1}^{\prime}\right\|_{2} 
+ M||R(W_1,\mathds{W},x_L)||_2||\mathds{W}-\mathds{W}^{\prime}||_2.
\end{gather*}
}

\begin{align*}
	\Delta_{L} & \leq\quad C_{\kw{COM}} B_x\|W_1-W_1^{\prime}\|_2 \\
	&\quad + 2C_{\kw{COM}} |R| B_2C_\rho C_{\kw{MSG}_0}C_{\kw{MSG}_1}d\max_{j\in C(x_L)}\Delta_{L-1,j} \\
	&\quad + C_{\kw{COM}}|R|\max_{r\in R}\|W_r-W_r^{\prime}\|_2\|R^+(W_1,\mathds{W},x_L,r)\|_2\\
	&\quad + C_{\kw{COM}}|R|\max_{r\in R}\|W_r-W_r^{\prime}||_2\|R^-(W_1,\mathds{W},x_L,r)\|_2\\
% & \leq\quad C_{\kw{COM}} B_x\|W_1-W_1^{\prime}\|_2 \\
%&\quad + 2C_{\kw{COM}} |R| B_2C_\rho C_{\kw{MSG}_0}C_{\kw{MSG}_1}d\max_{j\in C(x_L)}\Delta_{L-1,j} \\
%&\quad + 2 C_{\kw{COM}} |R| {\mathcal C} B_1  B_xM \max_{r\in R}\|W_r-W_r^{\prime}||_2\\
	&\leq\quad 4M \bigg(2B_{x}||W_{1}-W_{1}^{\prime}||_{2}\\
	&\quad+\max_{r\in R}\|W_r-
	W_r^{\prime}\|_2(\|R^+(W_1,\mathds{W},x_L,r)\|_2\\
	&\quad+\|R^-(W_1,\mathds{W},x_L,r)\|_2)\bigg) \tag{18}
\end{align*}



\eat{
\begin{equation}\Delta_L\leq MB_x\left|\left|W_1-W_1^{\prime}\right|\right|
_2+M\overline{R}||
W_r-W_r^{\prime}||_2\end{equation}
}

\stitle{$\bullet$ Step 3: Bound the prediction changes}.
We first reform the change of 
the prediction function as follows. 

\begin{equation*}\begin{aligned}
		P_{L} & \triangleq\quad\|T^Q_{L}(W_{1},\mathds{W})-T^G_{L}(W_{1},
		\mathds{W})\|_c \\
		& =\quad\bigg\|\kw{COM}(W_{1}x^Q_{L}+\sum_{r\in R} W_rR^+(W_{1},\mathds{W},x^Q_{L},r)), W_{1}x^Q_{L}\\
		& \quad\quad+\sum_{r\in R} W_rR^-(W_{1},\mathds{W},x^Q_{L},r))\\
			& \quad\quad-\kw{COM}(W_{1}x^G_{L}+\sum_{r\in R} W_rR^+(W_{1},\mathds{W},x^G_{L},r)), W_{1}x^G_{L}\\
			&\quad\quad+\sum_{r\in R} W_rR^-(W_{1},\mathds{W},x^G_{L},r)) \bigg\|_c\\
%		&=W_{1}(x^Q_{L}-x^G_{L}) + \sum_{r\in R} W_r(R(W_{1},
%		\mathds{W},x^Q_{L},r)-R(W_{1},\mathds{W},x^G_{L},r))\bigg\|_c\\
\end{aligned}\end{equation*}

Then the change of prediction can be bounded as follows.

\begin{scriptsize}
\begin{align*}
		\nabla_{L} & \triangleq \quad |P_L - P'_L|\\
		& =\quad\bigg|\|T^Q_{L}(W_{1},\mathds{W})-T^G_{L}(W_{1},
		\mathds{W})\|_c - \|T^Q_{L}(W'_{1},\mathds{W}')-T^G_{L}(W'_{1},
		\mathds{W}')\|_c\bigg| \\
		& \leq \quad\bigg|\|T^Q_{L}(W_{1},\mathds{W})-T^Q_{L}(W'_{1},
		\mathds{W}')-  (T^G_{L}(W_{1},
		\mathds{W})-T^G_{L}(W'_{1},
\mathds{W}'))\|_c\bigg| \tag{19}\\
		& \leq \quad\|T^Q_{L}(W_{1},\mathds{W})-T^Q_{L}(W'_{1},
\mathds{W}')\|_2 +
\|T^G_{L}(W_{1},
\mathds{W})-T^G_{L}(W'_{1},
\mathds{W}')\|_2 \\
&\leq \quad 4M \bigg(2B_{x}||W_{1}-W_{1}^{\prime}||_{2}+M\max_{r\in R}\|W_r-
W_r^{\prime}\|_2(\|R^+(W_1,\mathds{W},x_L,r)\|_2\\
&\quad\quad+\|R^-(W_1,\mathds{W},x_L,r)\|_2)\bigg)\tag{20}
\end{align*}
\end{scriptsize}


Equation (19) holds, due to the triangle inequality 
of norm $\|\cdot\|_c$.
And Equation (20) holds due to Equation (18) above.

\eat{
Here, we show that $\|a\|_c - \|b\|_c\leq \|a-b\|_c$.
Assume that $\|a\|_c = a_i$ and $\|b\|_c=b_j$.
Then $b_i\leq b_j$ and $-b_i\geq -b_j$. 
Hence, $a_i-b_j\leq a_i-b_i \leq \max_{k\in [1, d]} a_i-b_i$.
That is, $\|a\|_c - \|b\|_c\leq \|a-b\|_c$.
}

\stitle{$\bullet$ Step 4: Bound the covering numbers}.
To bound the change in $\nabla_{L}$ by 
a given threshold $\epsilon$, we use the covering number.
Intuitively, given a set $\M$ of matrices
that have a bound for the norm $\|\cdot\|_F$,
\ie $\|M\|_F\leq \lambda$ for all $M\in \M$,
and a number $\eta$, its covering 
number ${\mathcal N}(\M, \eta, \|\cdot\|_F)$ is defined as the minimum number
of balls of radius $\eta$, needed to cover
all matrices in $\M$~\cite{Covering-number}.
To bound the covering number, 
we use a lemma from~\cite{Covering-1}, 
which states that 
$${\mathcal N}(\M, \eta, \|\cdot\|_F)\leq \bigg(1+\frac{2d\lambda}{\epsilon}\bigg)^{d^2}$$

Here, $\varepsilon$ is a given threshold, $\lambda$
is the bound for the norm of matrices in $\M$, 
and $d$ is the dimension 
of the matrix.

Since the bound in Equation (20) has two kinds
of learned matrices, \ie $W_1$ and matrices in $\mathds{W}$.
Using this lemma, we deduce that 


\begin{align*}\mathcal{N}\left(W_{1},\frac{\epsilon}{8MB_{x}},\|\cdot\|_{F}\right)
\leq\left(1+\frac{16MB_{x}B_{1}\sqrt{d}}{\epsilon}\right)^{d^{2}}\end{align*}

\begin{align*}\mathcal{N}\left(W_r,\frac{\epsilon}{8M{\mathcal R}},
	\|\cdot\|_{F}\right)\leq\left(1+\frac{16M{\mathcal R}B_{2}\sqrt{d}}{\epsilon}
	\right)^{d^{2}} \quad (r\in R)
\end{align*}

Here, ${\mathcal R} = \max(\|R^+(W_1,\mathds{W},x_L,r)\|_2,
\|R^-(W_1,\mathds{W},x_L,r)\|_2) \leq {\mathcal C} B_1  B_xM$.
Therefore, the overall covering number for \HGIN is bounded
by the production of all these covering numbers~\cite{Gen-bound}.
\begin{align*}P=\left(1+\frac{16M\sqrt d\max\{B_xB_1,{\mathcal R}B_2\}}
{\epsilon}\right)^{2|R|d^2}\end{align*}

Because $\varepsilon$ is sufficiently small, 
\eg $\epsilon<4M\sqrt{d}\max\{B_xB_1,{\mathcal R}B_2\}$,
we have that $\log P\leq 2|R|d^2\log\left(\frac{32M\sqrt{d}\max\{B_xB_1,{\mathcal R}
B_2\}}{\epsilon}\right)$.
We will use this upper bound to prove the Rademacher complexity.

\eat{
\begin{equation*}
	3r^2\log\left(1+\frac{6M\sqrt{r}\max\{B_xB_1,\overline{R}B_2\}}{\epsilon}\right)
\end{equation*}
\begin{equation*}
	\epsilon<4 M\sqrt{r}\max\{B_xB_1,\overline{R}B_2\}
\end{equation*}
\begin{equation*}
	3r^2\log\left(\frac{8M\sqrt r\max\{B_xB_1,\overline{R}B_2\}}{\epsilon}\right)
\end{equation*}
}

\stitle{$\bullet$ Step 5: Bound for the Rademacher complexity}.
We finally can bound the Rademacher complexity.
To this end, we also apply a lemma from~\cite{RB-1},
which states that 
given a function ${\mathcal F}$ taking values
in $[0,1]$, 
its Rademacher complexity 
$\hat{\mathcal{R}}$ 
is bounded by $\inf_{\alpha>0}\left(\frac{4\alpha}{\sqrt{N_T}}+\frac{12}{N_T}\int_{\alpha}^{\sqrt{N_T}}\sqrt{\log\mathcal{N}(\mathcal{F},\epsilon,\|\cdot\|_F)}d\epsilon\right)$,
where $\mathcal{N}(\mathcal{F},\epsilon,\|\cdot\|_F)$ is the covering 
number of $\mathcal{F}$, and
$N_T$ is the number of training data.


From Step-4, we know that the covering number ${\mathcal N}$
is bounded by $\left(1+\frac{16M\sqrt d\max\{B_xB_1,{\mathcal R}B_2\}}
{\epsilon}\right)^{2|R|d^2}$.

Therefore, we can bound the covering number as follows.

\begin{align*}
\log \mathcal{N}(\mathcal{F},\epsilon,\|\cdot\|_F)&\leq 
\log \mathcal{N}(P,\frac{\epsilon}{2},\|\cdot\|_F)\\
&
\leq 2|R|d^2\log\left(\frac{32M\sqrt{d}\max\{B_xB_1,\overline{R}
	B_2\}}{\epsilon}\right)
\end{align*}

Setting $\alpha  = \frac{2}{\sqrt{N_T}}$, we have that 
$\hat{\mathcal{R}}(\mathcal{T})\leq\frac{8}{N_T}+
\frac{12}{\sqrt{N_T}}\sqrt{2|R|d^2\log T}$,
where $T=16M\sqrt{{N_T}d}\max\{B_xB_1,{\mathcal R}B_2\}$,
and $\mathcal{I}$ is the class of functions
mapping the computation of embedding to the prediction function $\nabla_{L}$.

We can verify that the max margin loss function 
is 1-Lipschitz.
Then ${\hat{\mathcal{R}}(\mathcal{J}_\gamma)}
%\leq \hat{\mathcal{R}}
%_T(\mathcal{I})
\leq\frac{8}{N_T}+\frac{12d}{\sqrt{N_T}}
\sqrt{2|R|d^2|R|\log T}$.

\stitle{The 1-Lipschitz of the loss function}.
We next show that the max margin loss function 
is 1-Lipschitz, 
\ie $|{\mathcal L}(h_u,h_v)-{\mathcal L}(h'_u,h'_v)|\leq 
|{\mathcal M}(h_u,h_v)-{\mathcal M}(h'_u,h'_v)|$.
Since the loss functions for positive data and negative data
are different,
we consider all possible four cases:
(1) both $(h_u,h_v)$ and $(h'_u,h'_v)$ are positive;
(2) both $(h_u,h_v)$ and $(h'_u,h'_v)$ are negative;
(3) $(h_u,h_v)$ is positive and $(h'_u,h'_v)$ is negative; and
(4) $(h_u,h_v)$ is negative and $(h'_u,h'_v)$ are positive.



\stab
(1) If both $(h_u,h_v)$ and $(h'_u,h'_v)$ are positive, 
then their loss are $-\alpha$.
Hence, 
$|{\mathcal L}(h_u,h_v)-{\mathcal L}(h'_u,h'_v)| = 0\leq 
|{\mathcal M}(h_u,h_v)-{\mathcal M}(h'_u,h'_v)|$.

\eat{
\stab
(2) If both $(Q_1,G_1)$ and $(Q_2,G_2)$ are negative,
and the margins between $(Q_1,G_1)$ and $(Q_2,G_2)$ are smaller than
$\alpha$, then the inequality holds.

\stab
(3) If both $(Q_1,G_1)$ and $(Q_2,G_2)$ are negative,
and the margins between $(Q_1,G_1)$ and $(Q_2,G_2)$ are larger than 
$\alpha$, then the loss is 0, and the inequality holds.
}

\stab
(2) If both $(h_u,h_v)$ and $(h'_u,h'_v)$ are negative,
when both ${\mathcal M}(h_u,h_v)$
and ${\mathcal M}(h'_u,h'_v)$ are not larger than 
$\alpha$, we have that 

\begin{align*}
&{\mathcal L}(h_u,h_v)-{\mathcal L}(h'_u,h'_v)|\\
&=(\alpha -\|\max(0,h_u-h_v)\|^2_2) - (\alpha - \|\max(0,h'_u-h'_v)\|^2_2)\\
&=-\|h_u-h_v\|^2_2 + \|h'_u-h'_v\|^2_2 \\
&\leq |{\mathcal M}(h_u,h_v)-{\mathcal M}(h'_u,h'_v)| \tag{21}
\end{align*}

Equation (21) holds, since 
$\mathcal{M}(h_u, h_v) = \cdd{\max(0, h_u-h_v)}_2^2$.

\vspace{2ex}
When ${\mathcal M}(h_u,h_v)\geq \alpha$ or
 ${\mathcal M}(h'_u,h'_v)\geq  \alpha$,
 the proof is similar.











\stab
(3) If $(h_u,h_v)$ is positive and $(h'_u,h'_v)$ is negative,
and ${\mathcal M}(h'_u,h'_v)\leq \alpha$, we have the following

$\|\max(0,h_u - h_v)\|_2 - \alpha  - (\alpha - \|\max(0,h'_u - h'_v)\|)\\
=   \|h'_u - h'_v \|_2- 2\alpha\\
\leq \|h_u - h_v\|_2 +  \|h'_u - h'_v\|_2- 2\|h'_u - h'_v \|_2\\
\leq \|h_u - h_v\|_2 - \|h'_u - h'_v \|_2\\
\leq |{\mathcal M}(h_u,h_v)-{\mathcal M}(h'_u,h'_v)|$





\stab
(4) If $(h_u,h_v)$ is positive and $(h'_u,h'_v)$ is negative,
and ${\mathcal M}(h'_u,h'_v)\geq \alpha$, then 
we have 

$
\|\max(0,h_u - h_v)\|_2 - \alpha  - \max(0,\alpha - \|\max(0,h'_u - h'_v)\|)\\
= - \alpha\\
\leq |{\mathcal M}(h_u,h_v)-{\mathcal M}(h'_u,h'_v)|$

\subsection*{A5. Complex queries}

\begin{figure}[tb!]
	%	\vspace{-6.7ex}
	\centerline{\includegraphics[width=0.8\columnwidth]{fig/query-example-2.eps}}
	\centering
	\vspace{-1.4ex}
	\caption{Example of complex query}
	\label{fig:query}
	\vspace{-1.4ex}
\end{figure}

When patterns are complex (\eg with multiple cycles
and large number of vertices),
the exact algorithms can be much slow 
and \HFrame is significantly faster than 
them. 
Consider a complex query extracted from Citeseerx,
which is depicted in 
Figure~\ref{fig:query}.
\kw{SIM}-\kw{TD}, \kw{EJ} and \kw{PJ} take
2.691s, 11.164s and 13.472s respectively to conduct the query,
while \HFrame takes only 0.081s.

\reviserdy{
\subsection*{A6. The Application of \HFrame}
We focus on the decision problem of the subgraph homomorphism problem, 
which can be applied in frequent pattern mining, graph pattern depenedency discovery, 
graph query evaluations, and accelerating subgraph enumerations. 

We can extend our algorithms to generate exact subgraph instances in the following two ways:
(1) plugging in the enumeration algorithm: 
we can apply \HFrame to verify the candidates of the first pattern vertex in the matching order, 
and apply exact algorithm to enumerate on the vertices, on which \HFrame returns true;
or (2) applying \HFrame to filter candidates: 
we can apply \HFrame to check whether there exists a match for each vertex v in C(u) for each pattern vertex u, 
remove all vertices in C(u) that does not have a match, 
and finally run the exact enumeration algorithm on the subgraph deduced by the remained vertices. 
Taking PJ as an example, we adopt the second method, 
and reports the results on Citation and DBLP as \ref{tab:app-filter}. 
We can see that \HFrame trades off a marginal loss in accuracy for a significant gain in computational efficiency.
\begin{table}[tb]
	\caption{Effectiveness of \HFrame in candidate filtering}
	\label{tab:app-filter}
	\vspace{-1.5ex}
	\centering
	\resizebox{1.01\columnwidth}{!}{
		\begin{tabular}{c||c|c|c|c} \hline
			& \multicolumn{2}{c|}{\bf{Accuracy}}
                & \multicolumn{2}{c}{\bf{Efficiency (s)}}\\
			\hline
			\diagbox{\bf{Method}}{\bf{Dataset}} &Citation&DBLP&Citation&DBLP\\\hline
			\kw{Only PJ}              & 1.000 & 1.000 & 1.545 & 1.803 \\ %\hline
			\kw{\HFrame+PJ}              & 0.815 & 0.832 & 0.562 & 0.547 \\ \hline
		\end{tabular}
	}
	\vspace{-2ex}
\end{table}

Additionally, in frequent pattern mining and graph pattern dependency discovery, 
the quality of patterns $Q$ is usually defined by the number of vertices matching the pivot of $Q$, 
rather than the number of matches. 
We have added an application of \HFrame in frequent pattern mining algorithm SPMiner\cite{ying2024representation}. 
We evaluated the performance on Enzyme dataset same as SPMiner, 
and calculated the median anchored frequency and grouped runtime by pattern size in \ref{tab:app-mining}. 
We found that when plugging \HFrame in the existing frequent pattern mining, 
\HFrame can discover high-quality patterns in a reasonable time.

\begin{table}[b]
        \caption{Effectiveness of \HFrame in frequent pattern mining}
            \label{tab:app-mining}
        \vspace{-1.5ex}
        \centering
        \resizebox{1.01\columnwidth}{!}{
            \begin{tabular}{c|c|c|c|c|c|c} \hline
                    \multicolumn{2}{c|}{} & \multicolumn{5}{c}{Pattern size} \\ \cline{3-7}
                    \multicolumn{2}{c|}{} & 5 & 6 & 7 & 8 & 9 \\ \hline
                    \multirow{2}{*}{Support} & SPMiner+\HFrame & 14,688 & 14,445 & 13,571 & 12,109 & 9,723 \\
                             & SPMiner & 13,005 & 14,572 & 13,691 & 9,831 & 7,928 \\ \hline
                    \multirow{2}{*}{Runtime(s)} & SPMiner+\HFrame & 160.85 & 193.77 & 226.01 & 256.43 & 292.16 \\
                             & SPMiner & 151.89 & 182.38 & 212.66 & 244.02 & 275.40 \\ \hline
            \end{tabular}
        }
        \vspace{-2ex}
\end{table}
}
	
\eat{

\clearpage
\subsection*{A5. Complex queries}

\begin{figure}[tb!]
	%	\vspace{-6.7ex}
	\centerline{\includegraphics[width=0.8\columnwidth]{fig/query-example-2.eps}}
	\centering
	\vspace{-1.4ex}
	\caption{Example of complex query}
	\label{fig:query}
	\vspace{-1.4ex}
\end{figure}

There exist some 

\warn{Figure~\ref{fig:query} illustrates a complex query example 
	that poses significant challenges for exact algorithms.}
}



\eat{
\begin{align*}
		\hat{\mathcal{R}}_{\mathcal{T}}(\mathcal{J}_{\gamma})\leq\frac{4}{\gamma m}+
		\frac{16r}{\gamma\sqrt{m}}\sqrt{3\operatorname{log}Q},where \\
		Q=16\sqrt{m}M\sqrt{r}\max\{B_{x}B_{1},\overline{R}B_{2}\}, \\
		M=C_\phi\frac{\left(\mathcal{C}d\right)^L-1}{\mathcal{C}d-1}, \\
		\overline{R}\leq C_{\rho}C_{g_0}C_{g_1}d\operatorname*{min}\left\{b\sqrt{r},B_{1}B_{x}
		M\right\}.
\end{align*}

\begin{align*}e\leq 1\end{align*}
\warn{Since we normalize the embeddings}

Hence, the empirical Rademacher complexity of $\mathcal{J}_\gamma$ with respect to $\mathcal{T}$ is
\begin{align*}\hat{\mathcal{R}}_{\mathcal{T}}(\mathcal{I})\leq\inf_{\alpha>0}\left(\frac{4\alpha}{\sqrt{m}}+\frac{12}{m}\int_{\alpha}^{2e\sqrt{m}}\sqrt{\log\mathcal{N}(\mathcal{I},\epsilon,\|\cdot\|)}d\epsilon\right)\end{align*}

\begin{align*}
		&\int_{\alpha}^{2e\sqrt{m}}\sqrt{\log\mathcal{N}(\mathcal{I},\epsilon,\|\cdot\|)}d\epsilon \\
		&\leq\int_{\alpha}^{2e\sqrt{m}}\sqrt{\log\mathcal{N}(\mathcal{B},\frac{\epsilon}{2},dist(\cdot,\cdot))}d\epsilon \\
		&\leq\int_{\alpha}^{2e\sqrt{m}}\sqrt{\log U}\leq2e\sqrt{m}\sqrt{\operatorname{log}U}
		\quad=2\sqrt{m\operatorname{log}U}
\end{align*}

\begin{align*}2|R|r^2\log\left(\frac{16M\sqrt r\max\{B_xB_1,\overline{R}B_2\}}{\alpha}
	\right)
\end{align*}

}